\section{Continuous inputs without a power-of-accuracy width cost}
\label{sec:quantized-input}
A continuous input can first be converted to binary digits.  Applying a
Boolean-network compiler to those digits avoids a separate neuron for every
small grid box.  The price is logarithmic additional depth.  This section
concerns only a comparison network used in a proof; the Bayesian learner
continues to use its raw output in both the likelihood and prediction.

\begin{lemma}[A tanh binary quantizer]\label{lem:tanh-quantizer}
Let $Q=2^t$, $t\ge1$, and $0<\rho<1/(4Q)$.  Let
$0<\delta\le\rho/[10^6Q(t+1)]$.  There is a tanh network of width at most
$d(t+2)+1$ and hidden depth $1+2t$ which, for every
$y\in[0,1]^d$ satisfying
\[
 \operatorname{dist}(y_j,Q^{-1}\mathbb Z)\ge\rho
 \quad(1\le j\le d),
\]
returns all $dt$ signed binary digits of $\lfloor Qy_j\rfloor$, each with
error at most $\delta$.  Its coefficients have absolute value at most
\[
 C\rho^{-1}\sqrt{Q(t+1)/\rho}\log(4/\delta).
\]
The output digits always lie in $[-1,1]$, including outside the displayed
set.  Arbitrary further layers can copy the digits while preserving their
error, with coefficient cap $C\log(4/\delta)$.
\end{lemma}
\begin{proof}
It suffices to process one coordinate in parallel with the others.  The
exact binary algorithm starts at $z_0=y$ and iterates
\[
 b_j=\mathbf 1\{z_j\ge1/2\},\qquad z_{j+1}=2z_j-b_j
 \quad(0\le j<t).
\]
Every threshold encountered is a dyadic grid boundary.  Therefore
$|z_j-1/2|\ge\rho$ under the stated hypothesis.  Choose
\[
 \tau^2=\frac{\rho}{10^6Q(t+1)},\qquad
 k=\frac4\rho\log(4/\delta),
\]
and initialize $u_0=\tanh(\tau y)$.  One physical layer computes
\[
 v_j=\tanh(u_j),\qquad
 s_j=\tanh\big(k(u_j/\tau-1/2)\big),
\]
and the next computes
\[
 u_{j+1}=\tanh\big(2v_j-\tfrac\tau2(1+s_j)\big).
\]
At each of these layers every previously stored signed bit is replaced by
$\tanh(k s)$; thus copies are restored rather than accumulated linearly.
Constants can be supplied by affine biases.  An exact constant wire of value
$\tanh(1)$ is an alternative if the model omits intermediate biases.

Write $e_j=|u_j/\tau-z_j|$.  As long as $e_j\le\rho/8$, the decision
input has margin at least $7\rho/8$, so
$|s_j-(2b_j-1)|\le\delta$.  The elementary bound
$|\tanh v-v|\le |v|^3/3$ gives
\[
 e_0\le\tau^2/3,\qquad
 e_{j+1}\le2e_j+100\tau^2+\delta/2.
\]
Here $|u_j|\le2\tau$ and the next preactivation has magnitude at most
$5\tau$, so the numerical constant $100$ is sufficient.  Consequently
$e_j\le102Q\tau^2+Q\delta/2<\rho/8$ for every $j\le t$;
induction justifies the temporary hypothesis.  A stored bit has magnitude
at least $1-\delta\ge7/8$, so its copy has error at most
$2e^{-7k/4}\le\delta$.  The displayed width, depth, and coefficient cap
follow from these formulas.  The same restoring copy proves the final claim.
\end{proof}

\begin{lemma}[Quantized-input raw tanh compilation]
\label{lem:quantized-raw-compiler}
There is a numerical exponent $b\ge2$ with the following property.  Let
$g$ be an arbitrary-real-coefficient ReLU network of hidden width $C\ge1$
and depth $h\ge1$.  Suppose $P_X$ has density at most $\kappa$ relative to the uniform
probability measure on $[-1,1]^d$, and put $\bar\kappa=\max\{1,\kappa\}$.  For
$0<\epsilon<1/4$, define
\[
 J=\left\lceil\log_2\frac{c\bar\kappa d(C+1)^h}{\epsilon}\right\rceil,
 \qquad
 K=4(dJ+1)(C+1)^h+\lceil\log_2(32/\epsilon)\rceil+2.
\]
At every prescribed depth $H\ge2J+b(h+1)+3$, there is a tanh network $q$
with raw affine readout (using the native output bias) such that
\[
 \E|q(X)-\operatorname{clip}(g(X),[-1,1])|^2\le c\epsilon,
 \qquad \sup_{x\in[-1,1]^d}|q(x)|\le4,
\]
and
\[
 m\le c K^b,\qquad
 \log(B+2)\le c\big[J+\log(d\bar\kappa/\epsilon)+\log(K+2)\big].
\]
The raw readout has absolute coefficient sum plus absolute bias at most
four.  These resource bounds are independent of additional learner depth.
The constants do not depend on any coefficient of $g$.
\end{lemma}
\begin{proof}
Write $f=\operatorname{clip}(g,[-1,1])$ and $R=(C+1)^h$.
On every coordinate line, a depth-$h$, width-$C$ ReLU network has at most
$R$ affine intervals.  Clipping each affine interval is monotone, so the
one-dimensional total variation of $f$ on any segment is at most $2R$.

Take $Q=2^J$ and mesh $\Delta=4/Q$.  For a grid shift
$a\in[0,\Delta)^d$, let
\[
 y=(x-a+2\mathbf 1)/4,\qquad
 x_a=a-2\mathbf 1+4Q^{-1}\lfloor Qy\rfloor.
\]
For $J$ sufficiently large, $y\in[0,1]^d$ whenever
$x\in[-1,1]^d$.  Averaging over $a$ and using the coordinate translation
bound for functions of bounded variation gives
\[
 \mathbb E_a\int_{[-1,1]^d}|f(x)-f(x_a)|\,2^{-d}dx\le c dR/Q.
\]
For completeness, replace the coordinates one at a time.  In one
coordinate the randomized rounding displacement lies in $[0,\Delta]$;
Fubini and
$\int|u(t+s)-u(t)|dt\le |s|\operatorname{TV}(u)$
bound its average contribution by $2R\Delta$, times the volume of the
other coordinates in the slightly enlarged cube.  After the factor $2^{-d}$,
this volume contributes at most $(1+\Delta)^d\le e$, because the numerical
constant in $J$ ensures $Q\ge8d$.  Sum the $d$ bounds.
Multiplying by the density bound shows that some fixed shift $a$ has
\[
 \E|f(X)-f(X_a)|\le c\bar\kappa dR/Q\le c\epsilon.
\]
The chosen shift is part of the proof's comparison network, not an input
to the learner.

Set $\rho=\epsilon/(c\bar\kappa dQ)$, increasing the numerical constant
$c$ if necessary.
The union of the original-coordinate strips for
$\operatorname{dist}(y_j,Q^{-1}\mathbb Z)<\rho$ has probability at most
$c\bar\kappa dQ\rho\le\epsilon$.  On its complement use
Lemma~\ref{lem:tanh-quantizer}.  The exact bits describe $X_a$ as an affine
function of $dJ$ signed Boolean variables.  Composing this affine decoding
with the first layer of $g$ gives a width-$C$, depth-$h$ ReLU network on
that Boolean cube.  The coefficient-free majority-circuit compiler used
in Lemma~\ref{lem:raw-Boolean-tanh} has size $S\le K^b$, depth
at most $b(h+1)$, and signed-bit readout variation plus offset below four.
It approximates $f(X_a)$ to $\epsilon/4$.

Choose the quantizer bit tolerance to be
\[
 \delta=\min\left\{\frac{\rho}{10^6Q(J+1)},
                    \frac1{100(S+1)},\frac\epsilon{100}\right\}.
\]
The majority gates have integer margin at least one.  Replacing each
by tanh with gain $\log(4/\delta)$ restores bit error at most $\delta$,
because its incoming perturbation is at most $S\delta\le1/100$.
The raw readout therefore adds at most $4\delta$ error.  Insert as many
restoring signed-bit layers as necessary before the readout to attain the
prescribed $H$.  This does not increase the error or coefficient cap.

On the good quantizer set, $|q-f(X_a)|\le\epsilon/2$.  Everywhere,
$|q|\le4$, because the last hidden coordinates lie in $[-1,1]$ and the
readout has variation plus offset below four.  Thus its squared error
against $f$ is bounded by $25$ on the bad set.  On the good set use
$|f-f(X_a)|^2\le2|f-f(X_a)|$ and the preceding translation estimate.
This proves the claimed mean-square error.  The quantizer and Boolean
compiler widths are bounded by $cK^b$.  Their coefficient logarithms
are $O(J+\log(d\bar\kappa/\epsilon)+\log(K+2))$, directly from the
quantizer formulas and the majority-gate gains.  No inverse precision
appears as a power in the width.
\end{proof}
